\chapter{序章：二足歩行は本当にバランスが良いのか}\label{ch:intro}

\section{背景}

フィジカル AI は、ロボットのように実世界で身体を持って動く AI である。
その開発には高価なハードウェアが必要で、試行錯誤のたびに機体を壊すリスクや時間のコストがかかる。
そこで、安価なコンピュータシミュレーションで、どこまで実機の設計・制御・学習に使える知見が得られるかが重要になる。

本報告書が出発点とするのは、フィクションから生まれた概念である。
AMBAC（Active Mass Balance Auto Control、能動的質量移動による自動姿勢制御）は、
『機動戦士ガンダム』\cite{R1} の作品世界を支える SF 考証として、1981 年のムック本『ガンダムセンチュリー』で示されたとされる \cite{R42}。
人型の機動兵器が宇宙空間で手足を振り、その反作用で姿勢を変えるという考え方である。

これは単なる空想ではない。宇宙飛行士が手足の動きで体の向きを変えられることは古くから研究されており \cite{R8}、
落下する猫が空中で向きを変える現象も同じ原理で説明される \cite{R4,R7}。
実際に、宇宙飛行士の大西卓哉さんは国際宇宙ステーションで脚を使った姿勢制御に AMBAC の名を挙げ、ガンダムファンの間で話題になった \cite{R43,R44}。

\section{問い}

ロボットを題材にしたアニメやゲームでは、人型の機体が当たり前のように二本脚で立ち、手足を振って姿勢を整える。
一方で同じ作品世界には、脚のない機体について整備兵が言う「あんなの飾りです」という台詞もある \cite{R2}。
本報告書の中心となる問いは次のとおりである。

\begin{quote}
\textbf{二足歩行は本当にバランスが良いのか？　なぜ人は二本脚なのか？}
\end{quote}

\section{本報告書の答え（要約）}

\begin{itemize}
\item \textbf{静的には、二足歩行のバランスは悪い。}
  足裏の支持範囲が狭く重心が高いので、10 kg の荷物を腕の先で持つだけで、静止していても転倒する（第~\ref{ch:legs}章）。
  平地での安定性は戦車型の方が上である。
\item \textbf{それでも立てるのは、動かせる重い上体による能動的なバランスがあるから。}
  上体を前後に傾ける腰の戦略 \cite{R31} は、手足の代わりに上体を動かす地上版の AMBAC であり、これで同じ荷物を運べるようになる。
  段差を上るときも、上体を前に傾けることで、膝だけでは 18 cm だった限界が 60 cm まで広がる（第~\ref{ch:steps}章）。
\item \textbf{人間の大きさだからこそ、脚が活きる。}
  人間大の二本脚は戦車型の 2 倍以上の段差を越えるが、大型化すると二乗三乗則 \cite{R33} で逆転し、18 m 級の脚は 5 cm の段差しか越えられない。
\item \textbf{条件次第で優位は入れ替わる。}
  宇宙・平地・段差、人間大・巨大機で、どの形状が有利かが変わる。
  したがって、用途や目的に応じて形状を合わせる必要がある（パーツを組み合わせて機体を構成するゲームの「アセンブル」\cite{R3} のように）。
\end{itemize}

\section{本報告書の構成}

本報告書は、プロジェクトの第 1 段階「形状と、それが与える意味」の成果をまとめたものである。

\begin{description}
\item[第 I 部：AMBAC の物理（宇宙）] 角運動量保存から AMBAC の基本式を導き、幾何学的位相として理解する（第~\ref{ch:theory}章、第~\ref{ch:planar}章）。
  リアクションホイールとの関係（第~\ref{ch:rw}章）、トルクとスケーリング則（第~\ref{ch:dyn}章）、学習と最適制御（第~\ref{ch:ai}章）を扱う。
\item[第 II 部：下半身の形状（地上）] 腕の反作用による転倒・滑り（第~\ref{ch:tank}章）、「足は飾りか」（第~\ref{ch:legs}章）、段差（第~\ref{ch:steps}章）を扱う。
\item[第 III 部：上半身（肩と肩甲骨）] 3 次元の肩と回転の非可換性（第~\ref{ch:3d}章）、コッドマンのパラドックスと腕組み（第~\ref{ch:shoulder}章）を扱い、
  フィジカル AI に必要な可動域などのスペックとトレードオフを考察する。
\item[第 IV 部：結論] 問いへの答え、スペックとトレードオフの整理、第 2 段階の構想（第~\ref{ch:conclusion}章）。
\item[付録] 式の導出（付録~\ref{app:deriv}）、数値計算の方法と検証（付録~\ref{app:num}）、パラメータと再現手順（付録~\ref{app:param}）。
\end{description}

各章の見出しにある Sim01〜Sim10 は、リポジトリ \url{\repo} の \texttt{sims/} にあるシミュレーションスクリプトに対応する。
すべての数値と図は、付録~\ref{app:param} の手順で再現できる。

\section{前提とスコープ}

\begin{itemize}
\item \textbf{剛体の仮定}：機体は剛体リンクの集まりとして扱う。柔軟体・流体・ケーブル張力は扱わない。
\item \textbf{理想的なアクチュエータ}：多くの評価で、関節のトルク・速度を理想的に出せると仮定する。
  油圧・電動の動特性は第 2 段階の課題とする。
\item \textbf{環境}：宇宙（自由浮遊、角運動量が保存される）と地上（接触力・摩擦が働く）。重力は地球・火星・月を扱う。
\item \textbf{モデルの次元}：主に平面（2 次元）モデルで原理を確かめ、肩の運動など本質的に 3 次元の現象は 3 次元モデルで扱う。
\item \textbf{準静的な評価}：転倒・滑り・段差は、主にゆっくり動く場合の限界を求める。勢いを使う動的な動作は扱わない。
\end{itemize}
